\documentclass[11pt, a4paper]{book}
\usepackage{listings}
\usepackage[utf8]{inputenc}
\usepackage[english,spanish]{babel}
\RequirePackage{verbatim}
\usepackage{fancyhdr}
\usepackage{graphicx}
\usepackage{afterpage}
\usepackage{pdfpages}
\usepackage{longtable}

\usepackage[pdfborder={000}]{hyperref} %referencia

% ********************************************************************
% Re-usable information
% ********************************************************************
\newcommand{\myTitle}{Título del proyecto\xspace}
\newcommand{\myDegree}{Grado en Ingeniería Informática\xspace}
\newcommand{\myName}{Ángel Gómez Martín\xspace}
\newcommand{\myProf}{Juan Julián Merelo Guervos\xspace}
\newcommand{\myFaculty}{Escuela Técnica Superior de Ingenierías Informática y de Telecomunicación\xspace}
\newcommand{\myFacultyShort}{E.T.S. de Ingenierías Informática y de
Telecomunicación\xspace}
\newcommand{\myDepartment}{Departamento de Arquitectura y Tecnología de Computadores\xspace}
\newcommand{\myUni}{\protect{Universidad de Granada}\xspace}
\newcommand{\myLocation}{Granada\xspace}
\newcommand{\myTime}{\today\xspace}
\newcommand{\myVersion}{Version 0.1\xspace}


\usepackage{url}
\usepackage{colortbl,longtable}
%\makeindex
%\usepackage[style=long, cols=2,border=plain,toc=true,number=none]{glossary}
%\makeglossary

%\renewcommand{\indexname}{Índice alfabético}
%\renewcommand{\glossaryname}{Glosario}

\pagestyle{fancy}
\fancyhf{}
\fancyhead[LO]{\leftmark}
\fancyhead[RE]{\rightmark}
\fancyhead[RO,LE]{\textbf{\thepage}}
\renewcommand{\chaptermark}[1]{\markboth{\textbf{#1}}{}}
\renewcommand{\sectionmark}[1]{\markright{\textbf{\thesection. #1}}}

\setlength{\headheight}{1.5\headheight}

\newcommand{\HRule}{\rule{\linewidth}{0.5mm}}
%Definimos los tipos teorema, ejemplo y definición podremos usar estos tipos
%simplemente poniendo \begin{teorema} \end{teorema} ...
\newtheorem{teorema}{Teorema}[chapter]
\newtheorem{ejemplo}{Ejemplo}[chapter]
\newtheorem{definicion}{Definición}[chapter]

\definecolor{gray97}{gray}{.97}
\definecolor{gray75}{gray}{.75}
\definecolor{gray45}{gray}{.45}
\definecolor{gray30}{gray}{.94}

\lstset{ frame=Ltb,
     framerule=0.5pt,
     aboveskip=0.5cm,
     framextopmargin=3pt,
     framexbottommargin=3pt,
     framexleftmargin=0.1cm,
     framesep=0pt,
     rulesep=.4pt,
     backgroundcolor=\color{gray97},
     rulesepcolor=\color{black},
     %
     stringstyle=\ttfamily,
     showstringspaces = false,
     basicstyle=\scriptsize\ttfamily,
     commentstyle=\color{gray45},
     keywordstyle=\bfseries,
     %
     numbers=left,
     numbersep=6pt,
     numberstyle=\tiny,
     numberfirstline = false,
     breaklines=true,
   }
 
% minimizar fragmentado de listados
\lstnewenvironment{listing}[1][]
   {\lstset{#1}\pagebreak[0]}{\pagebreak[0]}

\lstdefinestyle{CodigoC}
   {
	basicstyle=\scriptsize,
	frame=single,
	language=C,
	numbers=left
   }
\lstdefinestyle{CodigoC++}
   {
	basicstyle=\small,
	frame=single,
	backgroundcolor=\color{gray30},
	language=C++,
	numbers=left
   }

 
\lstdefinestyle{Consola}
   {basicstyle=\scriptsize\bf\ttfamily,
    backgroundcolor=\color{gray30},
    frame=single,
    numbers=none
   }


\newcommand{\bigrule}{\titlerule[0.5mm]}

%Para conseguir que en las páginas en blanco no ponga cabecerass
\makeatletter
\def\clearpage{%
  \ifvmode
    \ifnum \@dbltopnum =\m@ne
      \ifdim \pagetotal <\topskip
        \hbox{}
      \fi
    \fi
  \fi
  \newpage
  \thispagestyle{empty}
  \write\m@ne{}
  \vbox{}
  \penalty -\@Mi
}
\makeatother

\begin{document}
\begin{titlepage}
 
 
\newlength{\centeroffset}
\setlength{\centeroffset}{-0.5\oddsidemargin}
\addtolength{\centeroffset}{0.5\evensidemargin}
\thispagestyle{empty}

\noindent\hspace*{\centeroffset}\begin{minipage}{\textwidth}

\centering
\includegraphics[width=0.9\textwidth]{imagenes/logo_ugr.jpg}\\[1.4cm]

\textsc{ \Large TRABAJO FIN DE GRADO\\[0.2cm]}
\textsc{GRADO EN INGENIERÍA INFORMÁTICA}\\[1cm]
% Upper part of the page
% 
% Title
{\Huge\bfseries IPManager\\}
\noindent\rule[-1ex]{\textwidth}{3pt}\\[3.5ex]
{\large\bfseries Infraestructura bajo demanda}
\end{minipage}

\vspace{2.5cm}
\noindent\hspace*{\centeroffset}\begin{minipage}{\textwidth}
\centering

\textbf{Autor}\\{Ángel Gómez Martín}\\[2.5ex]
\textbf{Director}\\Juan Julián Merelo Guervós\\[2cm]
\includegraphics[width=0.3\textwidth]{imagenes/etsiit_logo.png}\\[0.1cm]
\textsc{Escuela Técnica Superior de Ingenierías Informática y de Telecomunicación}\\
\textsc{---}\\
Granada, Julio de 2020
\end{minipage}
%\addtolength{\textwidth}{\centeroffset}
%\vspace{\stretch{2}}
\end{titlepage}



\begin{titlepage}
 
 
\setlength{\centeroffset}{-0.5\oddsidemargin}
\addtolength{\centeroffset}{0.5\evensidemargin}
\thispagestyle{empty}

\noindent\hspace*{\centeroffset}\begin{minipage}{\textwidth}

\centering


 \vspace{3.3cm}

%si el proyecto tiene logo poner aquí
\includegraphics{imagenes/logo.png} 
 \vspace{0.5cm}

% Title

{\Huge\bfseries Título del proyecto}
\noindent\rule[-1ex]{\textwidth}{3pt}\\[3.5ex]
{\large\bfseries Subtítulo del proyecto.\\[4cm]}
\end{minipage}

\vspace{2.5cm}
\noindent\hspace*{\centeroffset}\begin{minipage}{\textwidth}
\centering

\textbf{Autor}\\ {Ángel Gómez Martín}\\[2.5ex]
\textbf{Director}\\
{Juan Julián Merelo Guervos}\\[2cm]
%\includegraphics[width=0.15\textwidth]{imagenes/tstc.png}\\[0.1cm]
%\textsc{Departamento de Teoría de la Señal, Telemática y Comunicaciones}\\
%\textsc{---}\\
Granada, mes de 2020
\end{minipage}
%\addtolength{\textwidth}{\centeroffset}
\vspace{\stretch{2}}

 
\end{titlepage}




\cleardoublepage

\chapter*{}
\thispagestyle{empty}

\begin{center}
{\large\bfseries IPManager: Infraestructura bajo demanda}\\
\end{center}
\begin{center}
Ángel Gómez Martín
\end{center}

\noindent{\textbf{Palabras clave}: servicios, infraestructura, aprovisionamiento, backend, frontend, \textit{API REST}}\\

\vspace{0.7cm}
\noindent{\textbf{Resumen}}\\

El aprovisionamiento, despliegue de servicios y almacenamiento de configuraciones de máquinas a gran escala son tareas que tradicionalmente se han desarorollado de forma independiente. Esto supone a corto y largo plazo una gran inversión de tiempo para los administradores de sistemas.

\bigskip
En este proyecto se ha desarrollado una solución compuesta por un backend y un frontend. Esta busca unificar todos estos procesos en un puesto centralizado que permite agilizar el desarrollo de estas tareas. También ofrece modularidad en dos grandes aspectos: permitiendo que se puedan agregar nuevas funcionalidades de forma sencilla en el backend y proporcionando una \textit{API REST} que permite la comunicación con otro tipo de aplicaciones.


\cleardoublepage
\thispagestyle{empty}

\begin{center}
{\large\bfseries Project Title: Project Subtitle}\\
\end{center}
\begin{center}
Ángel Gómez Martín
\end{center}

\noindent{\textbf{Keywords}: services, infrastructure, provisioning, backend, frontend, \textit{REST API}}\\

\vspace{0.7cm}
\noindent{\textbf{Abstract}}\\

Provisioning, deployment of services and storage of large-scale machine configurations are tasks that have traditionally been carried out independently. In the short and long term, this is a major time investment for system administrators.

\bigskip
In this project, a solution composed of a backend and a frontend has been developed. It seeks to unify all these processes in a centralized position that allows to speed up the course of these tasks. It also offers modularity in two major aspects: allowing new functionalities to be added easily in the backend and providing a \textit{REST API} that allows communication with other type of applications.











\chapter*{}
\thispagestyle{empty}

\noindent\rule[-1ex]{\textwidth}{2pt}\\[4.5ex]

Yo, \textbf{Ángel Gómez Martín}, alumno de la titulación Grado en Ingeniería Informática de la \textbf{Escuela Técnica Superior de Ingenierías Informática y de Telecomunicación de la Universidad de Granada}, con DNI XXXXXXXXX, autorizo la ubicación de la siguiente copia de mi Trabajo Fin de Grado en la biblioteca del centro para que pueda ser consultada por las personas que lo deseen.

\vspace{6cm}

\noindent {Fdo: Ángel Gómez Martín}

\vspace{2cm}












\chapter*{}
\thispagestyle{empty}

\noindent\rule[-1ex]{\textwidth}{2pt}\\[4.5ex]

D. \textbf{Juan Julián Merelo Guervós}, Profesor del Área de XXXX del Departamento de Arquitectura y Tecnología de Computadores de la Universidad de Granada.

\vspace{0.5cm}

\textbf{Informa:}

\vspace{0.5cm}

Que el presente trabajo, titulado \textit{\textbf{IPManager, Infraestructura bajo demanda}},
ha sido realizado bajo su supervisión por \textbf{Ángel Gómez Martín}, y autoriza la defensa de dicho trabajo ante el tribunal que corresponda.

\vspace{0.5cm}

Y para que conste, expide y firma el presente informe en Granada a X de Julio de 2020.

\vspace{1cm}

\textbf{El director:}

\vspace{5cm}

\noindent {Fdo: Juan Julián Merelo Guervós}








\chapter*{Agradecimientos}
\thispagestyle{empty}
\vspace{1cm}

A mi tutor, por la ayuda, conocimientos e ideas prestadas para la realización de este proyecto.

\bigskip
A mis padres y mi hermana, por todo el cariño y apoyo que me han dado durante toda esta etapa.



%\frontmatter
%\tableofcontents
%\listoffigures
%\listoftables
%
%\mainmatter
%\setlength{\parskip}{5pt}

\chapter{Introducción}


Cada día el número de equipos que se utilizan de forma profesional en empresas y organizaciones crece de forma exponencial. Todos ellos requieren una configuración inicial para empezar a trabajar, la cual puede comprender tanto la instalación de aplicaciones de forma local como el despliegue de servicios para poder trabajar con ellos, entre otros.

Esta configuración inicial se suele realizar en varias etapas, comenzando por la instalación física de los equipos, configuración de red, instalación de software y finalmente la configuración de este. Cada uno de estos pasos comprende una serie de tareas secundarias que generalmente se suelen realizar de forma independiente. La independencia de las tareas hace que la puesta a punto de todas las máquinas tome un tiempo muy valioso. Un ejemplo sería la instalación de un mismo software en cientos de máquinas, en la que hacerlo de manera secuencial tomaría mucho tiempo. Para solucionar esto han surgido herramientas que permiten realizar estas instalaciones en un solo paso, algo que soluciona en parte el problema.

Por otro lado el despliegue de los servicios que pueden necesitarse también suele tomar un tiempo considerable, algo que aunque está ligado al paso anterior, no se realiza al mismo tiempo.

Además todas las configuraciones de las máquinas que se manejan no suelen estar centralizadas. Cierto es que los routers o switches actuales tienen capacidades para almacenar todos esos datos, pero implica que se tiene que acceder a ellos para consultarlos.

\smallskip
En este proyecto se desarrolla una solución al problema expuesto, \textbf{IPManager}. Esta se compone de un backend y un frontend y busca unificar todos los procesos descritos anteriormente en un puesto centralizado que permite a un administrador de sistemas agilizar el desarrollo de estas tareas. También ofrece modularidad en dos grandes aspectos: permitiendo que se puedan agregar nuevas funcionalidades de forma sencilla en el backend y proporcionando una \textit{API REST}.

Con la reciente crisis provocada por el COVID-19 muchas empresas han tenido que dejar de lado sus oficinas para operar mediante teletrabajo. Herramientas como \textbf{IPManager} permiten que el administrador del sistema pueda administrarlo todo de forma remota sin tener que desplazarse de su hogar. Todas las facilidades que se puedan brindar al trabajo remoto son necesarias, por lo que ofrecer soluciones que permitan esto tiene una especial importancia.

\section{Objetivos}

El desarrollo de este proyecto tiene los siguientes objetivos:

\begin{itemize}
\item Ofrecer una solución que unifique todos estos procesos, para simplificarlos y para reducir el tiempo empleado en ellos. Estos son el aprovisionado de sistemas, el despliegue de servicios y el almacenado de configuraciones de máquinas.
\item Ofrecer una solución liviana y no intrusiva en el ecosistema donde se instale.
\item Ofrecer una solución que provea una de \textit{API REST} para permitir el desarrollo de otros frontend y la integración de estos.
\item Ofrecer una solución totalmente compuesta por software libre.
\end{itemize}

\section{Estructura del documento}

En este documento se explica todo el desarrollo del proyecto, así como las decisiones tomadas y los diferentes aspectos que se han tenido en cuenta.

Se comienza con una descripción más concreta del problema y con el análisis de requisitos en el \textit{capítulo 2}. Tras esto, en el \textit{capítulo 3}, se hace un repaso al estado del arte en este ámbito.

El capítulo 4 se ha dedicado a la exposición de las herramientas y tecnologías que se han empleado, y en el \textit{capítulo 5} y \textit{capítulo 6} se desarrolla en profundidad la solución que se propone y la implementación de la misma.

Finalmente, en el \textit{capítulo 7}, se exponen algunas conclusiones personales tras desarrollar este proyecto y también se enumeran los posibles trabajos futuros que se podrían incluir en este proyecto.

De forma anexa se ha incluido también la especificación de los endpoints del backend y de las variables de entorno, una guía de instalación del proyecto y una introducción a los primeros pasos a dar tras instalar el software.
%
%\input{capitulos/02_EspecificacionRequisitos}
%
%\input{capitulos/03_Planificacion}
%
%\input{capitulos/04_Analisis}
%
%\input{capitulos/05_Diseno}
%
%\input{capitulos/06_Implementacion}
%
%\input{capitulos/07_Pruebas}
%
%\input{capitulos/08_Conclusiones}
%
%%\chapter{Conclusiones y Trabajos Futuros}
%
%
%%\nocite{*}
%\bibliography{bibliografia/bibliografia}\addcontentsline{toc}{chapter}{Bibliografía}
%\bibliographystyle{miunsrturl}
%
%\appendix
%\input{apendices/manual_usuario/manual_usuario}
%%\input{apendices/paper/paper}
%\input{glosario/entradas_glosario}
% \addcontentsline{toc}{chapter}{Glosario}
% \printglossary

\end{document}
